\documentclass[a4paper,11pt]{article}

% Pacotes Principais -----------------------------------------------------------
\usepackage[portuges,brazil]{babel}
\usepackage[utf8]{inputenc}

% Formatação de capítulos ------------------------------------------------------
%\usepackage[Sonny]{fncychap}
%\usepackage{fncychap}
\usepackage{capitulos}

% Figuras e Imagens ------------------------------------------------------------
\usepackage{graphicx}
% Figuras lado a lado
\usepackage{epsfig}
\usepackage{subfigure}

% Utilizar H para inserir as imagens REALMENTE onde eu desejo
\usepackage{float}

% Fontes -----------------------------------------------------------------------
\usepackage[T1]{fontenc}
\usepackage{pslatex}

% Simbolos ---------------------------------------------------------------------
\usepackage{textcomp}
\usepackage{bbding}

% Tabelas ----------------------------------------------------------------------
%\usepackage{multicol}
\usepackage{multirow}
% Colorir a tabela
\usepackage{colortbl}
% Pacote hhline corrige os bugs das linhas que não aparecem com o colortbm
\usepackage{hhline}
% Tabelas com colunas de largura auto ajustável
\usepackage{tabularx}
% Notas de rodapé em tabelas (Pode-se usar o ambiente longtable também -
% Pesquisar exemplo com longtable)
\usepackage{threeparttable}
% Tabelas grandes
\usepackage{supertabular}

% Glossário --------------------------------------------------------------------
\usepackage[portuguese,noprefix]{nomencl}
\usepackage{makeglo}

% Outros pacotes ---------------------------------------------------------------
\usepackage{noitemsep}

% Comentários em bloco
\usepackage{verbatim}

% Hiperlinks
\usepackage{hyperref}

% Orientação de página
\usepackage{lscape}

% utilitários matemáticos
\usepackage{amsmath}
\usepackage{icomma}

% Evita o problema "too many unprocessed floats", colocando os campos
% 'flutuantes' em suas respectivas seções
\usepackage[section]{placeins}

% Referências ------------------------------------------------------------------
\usepackage[abbr]{harvard}	% As chamadas são sempre abreviadas
\harvardparenthesis{square}	% Colchetes nas chamadas
\harvardyearparenthesis{round}	% Parêntesis nos anos das referências
\renewcommand{\harvardand}{e}	% Substituir "&" por "e" nas referências

% Comandos gerais --------------------------------------------------------------
\newcommand{\titulo}{Sistemas de Tempo Real}
\newcommand{\autor}{Diogo Leite Rebouças}

% Configuração da fonte
%\renewcommand{\familydefault}{\sfdefault}

% Margens ----------------------------------------------------------------------
\setlength{\oddsidemargin}{3.5cm}
\setlength{\evensidemargin}{2.5cm}
\setlength{\textwidth}{15cm}
\addtolength{\oddsidemargin}{-1in}
\addtolength{\evensidemargin}{-1in}

\setlength{\topmargin}{2.0cm}
\setlength{\headheight}{1.0cm}
\setlength{\headsep}{1.0cm}
\setlength{\textheight}{22.7cm}
\setlength{\footskip}{1.0cm}
\addtolength{\topmargin}{-1in}

% Glossário --------------------------------------------------------------------
\makeglossary

% Capítulos --------------------------------------------------------------------
% Não aparecer o número na primeira página dos capítulos
\newcommand{\mychapter}[1]{\chapter{#1}\thispagestyle{empty}}
\newcommand{\mychapterstar}[1]{\chapter*{#1}\thispagestyle{empty}}

% Comandos matemáticos ---------------------------------------------------------
% Implicação em fórmulas
\newcommand{\implica}{\quad\Rightarrow\quad} %Meio de linha
\newcommand{\implicafim}{\quad\Rightarrow}   %Fim de linha
\newcommand{\tende}{\rightarrow}

% Fração com parenteses
\newcommand{\pfrac}[2]{\parent{\frac{#1}{#2}}}

% Transformada de Laplace e transformada Z
\newcommand{\lapl}{\pounds}
\newcommand{\transfz}{\mathcal{Z}}

% Sequências
\newcommand{\sequencia}[4]{$#1_{#2}$, $#1_{#3}$, \ldots, $#1_{#4}$}

% Outros ----------------------------------------------------------------------
\newcommand{\chave}[1]{\left\{#1\right\}}
\newcommand{\colchete}[1]{\left[#1\right]}
\newcommand{\parent}[1]{\left(#1\right)}

\newcommand{\rhoagua}{\rho_{\tiny \text{\tiny H}_2\text{\tiny O}}}

\let\D\displaystyle
\newcommand{\reg}{\textsuperscript{\textregistered}}
\newcommand{\soft}{\textit{software}}
\newcommand{\Soft}{\textit{Software}}

% Imagens exportadas pelo gnuplot ----------------------------------------------
\graphicspath{{imgs/eps/}}


\begin{document}

\begin{table}[H]
\centering
\begin{tabular}{cl}
% Primeira Linha
\multirow{7}{*}
{
\includegraphics[width=1.65cm]{imgs/uern}
} & \\
& Governo do Estado do Rio Grande do Norte\\
% Segunda Linha
& Secretaria de Estado da Educação, da Cultural e dos Desportos -- SECD\\
% Terceira Linha
& {\sc Universidade do Estado do Rio Grande do Norte -- UERN}\\
% Quarta Linha
& Pró-Reitoria de Ensino e Graduação -- PROEG\\
% Quinta Linha
& Ciências da Computação -- Reposição da 1\textordfeminine\ Avaliação (10
  pontos)\\
& 
\end{tabular}
\end{table}

Aluno:

\vspace{0.25cm} 

Data: %\today

\section*{Questão 1 (2 pontos)}
Quando falamos sobre a representação da informação, mostramos que existem dois
tipos de representação de imagens. Cite quais são esses dois tipos de
representação e determine suas características e aplicações.

\pagebreak

\section*{Questão 2 (1 ponto)}
Durante o processo de conversão A/D, devemos determinar o tipo de quantização a
ser utilizada e o número de bits para se representar cada amostra do sinal.
Baseado nessas informações, responda as alternativas abaixo:

\begin{itemize}
    \item [a)] Cite e determine as características dos tipos de quantização
explicados em sala de aula.
    \item [b)] Determine qual será a velocidade mínima ({\bf em Kbit/s}) que um
canal deve assegurar para que haja uma transmissão sem perda de um sinal
digital, obtido a partir de um processo de digitalização (conversão A/D) de um
sinal de voz de 4,5 KHz. Cada amostra deve ser representada por 16 bits.
\end{itemize}

\pagebreak

\section*{Questão 3 (1 ponto)}
Sabemos que a criação de imagens, sons e vídeos no formato digital traz
benefícios até mesmo quando o produto final for distribuído em mídia analógica e
que esses recursos estão intimamente relacionados com o conceito de {\bf
multimídia}. Assim sendo, responda:

\begin{itemize}
    \item[a)] O que é multimídia?
    \item[b)] Quais são as principais características dos arquivos multimídia?
\end{itemize}

\pagebreak

\section*{Questão 4 (2 pontos)}
Mostramos em sala de aula que existem três tipos de produtos multimídia:
títulos, aplicativos e sítios. Baseado nessas informações, responda:

\begin{itemize}
    \item[a)] Quais são as principais diferenças entre os {\bf aplicativos com
interface multimídia} e os {\bf aplicativos multimídia}?
    \item[b)] O que são sítios? Quais são os possíveis materiais que constituem
um sítio?
\end{itemize}

\pagebreak

\section*{Questão 5 (2 pontos)}
Falamos em sala de aula que, para criar qualquer tipo de produto multimídia,
devemos considerar os aspectos relacionados aos {\bf direitos autorais} do
material que utilizamos. Sabendo disso, cite quais são os principais tipos de
restrições (limitações) a serem consideradas.

\pagebreak

\section*{Questão 6 (2 pontos)}
Uma imagem, representada no {\bf formato matricial}, semelhante à Fig.
\ref{fig:matriz}, utiliza uma codificação RGBA de 8 bits por cor ({\bf 8
bits/cor}) e {\bf 8 bits de transparência} para cada um de seus {\it pixels}.
Assim, responda:

\begin{figure}[H]
\centering
\includegraphics[width=0.3\textwidth]{imgs/matriz}
\caption{Matriz de {\it pixels} utilizada na Questão 6.}
\label{fig:matriz}
\end{figure}

\begin{itemize}
    \item [a)] Qual será o tamanho de um arquivo ({\bf em KBytes}) dessa imagem
sabendo que esta possui 60 linhas ($M = 60$) $\times$ 90 colunas ($N = 90$)?
    \item [b)] Deseja-se comprimir esta imagem sem que haja perda da informação.
Mas, para isso a imagem deve ser antes dividida em partes menores de $6 \times
9$ {\it pixels}. Em quantas partes essa imagem será decomposta? Qual será o
tamanho ({\bf em Bytes}) de cada uma dessas partes antes de passar pelo
algoritmo de compressão?
    \item [c)] Se cada parte da alternativa {\it b} sofrer uma compressão de
20\%, qual será o tamanho final ({\bf em KBytes}) do arquivo comprimido?

\end{itemize}

\pagebreak

\end{document}
